% P05.1 one-dimensional literal regularity classification proof, 26 September 2026.

\begin{proof}
By
\Cref{thm:frontier-polyhedral-zero-degree-integral-regularity},
$\deg(\gamma)=0$ is equivalent to compactness and implies literal
regularity.  It remains to treat positive degree.

Choose the point $v$ in a noncompact component and the rational number $C$
as in the statement.  The maximum of two polyhedrally convex functions is
polyhedrally convex.  Moreover, positivity at infinity makes $\gamma$ tend
to $+\infty$ along every unbounded end.  Hence
\[
 \varphi=\max\{\gamma,C\}\in
 \mathrm{PAPC}_{\mathbf Z}(\mathbf X,\gamma),
 \qquad
 \varphi-\gamma\in\mathrm{PA}_{\mathbf Z,b}(\mathbf X),
\]
and $\varphi=\gamma$ outside a compact subset of every noncompact
component.  The crossing points of $\gamma$ with the rational constant
$C$ are rational because the slopes of $\gamma$ are integral and its
affine constants are rational.  We may therefore choose a finite rational
subdivision $\Pi$ on which both functions are affine and on whose unbounded
cells they agree.

The inequality $C>\gamma(v)$ and equality near infinity show that
$\varphi-\gamma$ is nonconstant.  It consequently has a nonzero integral
slope $k_e$ on some oriented bounded edge cell $e$ of $\Pi$.  Fix an
integer $N>|k_e|$ and set
\[
 d:=\frac{\varphi-\gamma}{N},
 \qquad
 P:=\gamma+d
   =\left(1-\frac1N\right)\gamma+\frac1N\varphi.
\]
The function $P$ is polyhedrally convex and differs boundedly from
$\gamma$.

We construct a literal obstacle above $d$ which has the same values at
every vertex of $\Pi$.  Orient each bounded edge cell $I$, denote its
integral-affine length by $L_I$, and write $r_I\in\mathbf Q$ for the slope
of $d$ there.  If $r_I$ is an integer, set $u_N=d$ on $I$.  Otherwise put
\[
 m_I:=\lfloor r_I\rfloor,
 \qquad
 \alpha_I:=r_I-m_I\in\mathbf Q\cap(0,1),
\]
and define $u_N$ on $I$ to have slope $m_I+1$ on the initial segment of
length $\alpha_IL_I$ and slope $m_I$ on the remaining segment, with the
same initial value as $d$.  The average slope is $r_I$, so $u_N$ and $d$
have the same value at both endpoints.  The slope of $u_N$ decreases once;
hence $u_N$ is concave on $I$ and lies above the affine chord $d|_I$.
On every unbounded cell put $u_N=d=0$.  The endpoint values glue, all new
breakpoints and constants are rational, and every slope is integral.
Thus
\begin{equation}
 u_N\in\mathrm{PA}_{\mathbf Z,b}(\mathbf X),
 \qquad
 u_N\geq d,
 \qquad
 u_N=d\quad\text{on }\Pi^{(0)}.
 \label{eq:frontier-polyhedral-literal-roof}
\end{equation}

It follows from
\eqref{eq:frontier-polyhedral-literal-roof} that
$P\leq\gamma+u_N$, so $P$ is an admissible competitor for
$P_\gamma(u_N)$.  Conversely, let $h\leq\gamma+u_N$ be any admissible
polyhedrally convex competitor.  At every vertex $x$ of $\Pi$,
\[
                         h(x)\leq\gamma(x)+d(x)=P(x).
\]
On a bounded edge cell, ordinary one-variable convexity bounds $h$ by the
affine chord joining these endpoint bounds; that chord is $P$ because $P$
is affine on the cell.  On every unbounded cell the obstacle itself equals
$P$.  Hence $h\leq P$ on $\mathbf X$, proving
\eqref{eq:frontier-polyhedral-literal-classification-envelope}.

On the chosen cell $e$, the slope of $P$ is
\[
                         q_e=s_e+\frac{k_e}{N},
 \qquad s_e\in\mathbf Z.
\]
Since $0<|k_e|<N$, one has $q_e\notin\mathbf Z$.  Let $x_0$, $x_1$, and
$x_2$ divide $e$ into two subintervals of equal integral-affine length
$L_e/2$.  For any literal integral-slope polyhedrally convex function
$\psi$, its left derivative $m$ at $x_1$ is an integer.  If
$m\leq\lfloor q_e\rfloor$, convexity gives
\[
 \psi(x_1)-\psi(x_0)\leq m\frac{L_e}{2},
 \qquad
 P(x_1)-P(x_0)=q_e\frac{L_e}{2}.
\]
If $m\geq\lceil q_e\rceil$, monotonicity of the one-sided derivatives gives
\[
 \psi(x_2)-\psi(x_1)\geq m\frac{L_e}{2},
 \qquad
 P(x_2)-P(x_1)=q_e\frac{L_e}{2}.
\]
The two cases exhaust the integral values of $m$.  In either case, the
difference of two pointwise errors has absolute value at least
\[
              \frac{L_e}{2}\operatorname{dist}(q_e,\mathbf Z).
\]
At least one of those errors therefore has absolute value at least half
this quantity.  This proves
\eqref{eq:frontier-polyhedral-literal-classification-separation}.

The envelope in
\eqref{eq:frontier-polyhedral-literal-classification-envelope} cannot be a
uniform limit of literal integral-slope polyhedrally convex functions, so
it is not literally regularizable.  Thus positive degree implies failure
of literal regularity.  Together with the zero-degree theorem, this proves
all three equivalences.
\end{proof}
